\begin{figure}[tb]
\centering
\begin{tikzpicture}[
  x=1cm,y=1cm,font=\small,
  pathline/.style={draw=figRotation,line width=1.1pt},
  seamline/.style={draw=figSeam,line width=.8pt},
  seammap/.style={draw=figSeam,line width=.8pt,
    dashed,rounded corners=6pt,-{Stealth[length=1.7mm]}},
  guide/.style={draw=black!21,line width=.5pt},
  arrow/.style={-{Stealth[length=1.8mm]},line width=.8pt}
]
  \path[use as bounding box] (0,-.15) rectangle (13.95,5.13);
  \node[anchor=west,font=\small\bfseries] at (.02,4.92)
    {(a) Periodic body frame};
  \node[anchor=west,font=\small\bfseries] at (7.12,4.92)
    {(b) Parallel frame};

  % Laboratory-orientation dials. They are identical in both panels.
  \foreach \x/\a/\word in {.90/0/start,6.05/-90/end,8.00/0/start,13.15/-90/end} {
    \draw[guide] (\x,4.15) circle[radius=.23];
    \draw[guide,densely dotted] (\x,4.15)--({\x+.20},4.15);
    \draw[pathline,arrow] (\x,4.15)--({\x+.205*cos(\a)},{4.15+.205*sin(\a)});
    \node[font=\scriptsize,text=black!60] at (\x,4.54) {\word};
  }
  \node at (3.47,4.15) {$\mathcal A=\rho\,d\tau$};
  \node at (10.57,4.15) {$\widetilde{\mathcal A}_{\tau}=0$};

  % Periodic body section. Blue horizontal transport has coordinate
  % theta=-rho*tau. Equal heights on the two walls are identified.
  \fill[figInk!3] (.90,1.50) rectangle (6.05,3.50);
  \draw[guide] (.90,1.50)--(6.05,1.50);
  \draw[guide] (.90,3.50)--(6.05,3.50);
  \draw[seamline] (.90,1.50)--(.90,3.50);
  \draw[seamline] (6.05,1.50)--(6.05,3.50);
  \draw[guide,densely dashed] (.90,3.15)--(6.05,3.15);
  \draw[guide,densely dashed] (.90,1.80)--(6.05,1.80);
  \node[anchor=east,inner sep=3pt,font=\scriptsize] at (.90,3.15) {$0$};
  \node[anchor=west,inner sep=3pt,font=\scriptsize] at (1.00,1.85) {$-\pi/2$};
  \node[font=\scriptsize] at (.90,3.73) {$\theta$};

  \draw[seammap]
    (6.05,1.80)--(6.44,1.80)--(6.44,.72)--(.51,.72)--(.51,1.80)--(.90,1.80);
  \node[fill=white,inner sep=3pt,font=\scriptsize,text=figSeam]
    at (3.47,.72) {identity seam};
  \draw[seamline,fill=white] (.90,1.80) circle[radius=2.0pt];
  \draw[pathline,arrow] (.90,3.15)--(6.05,1.80);
  \fill[figRotation] (.90,3.15) circle[radius=1.7pt];
  \fill[figRotation] (6.05,1.80) circle[radius=1.7pt];
  \node[fill=figInk!3,inner sep=2pt,text=figRotation] at (3.58,2.86)
    {$\theta=-\rho\tau$};

  % Mechanical-parallel section. The same transport now has constant
  % coordinate theta_tilde=0. Its endpoint returns to theta_tilde=-pi/2
  % when expressed in the initial section across the rotated seam.
  \fill[figInk!3] (8.00,1.50) rectangle (13.15,3.50);
  \draw[guide] (8.00,1.50)--(13.15,1.50);
  \draw[guide] (8.00,3.50)--(13.15,3.50);
  \draw[seamline] (8.00,1.50)--(8.00,3.50);
  \draw[seamline] (13.15,1.50)--(13.15,3.50);
  \draw[guide,densely dashed] (8.00,1.80)--(13.15,1.80);
  \node[anchor=east,inner sep=3pt,font=\scriptsize] at (8.00,3.15) {$0$};
  \node[anchor=west,inner sep=3pt,font=\scriptsize] at (8.10,1.85) {$-\pi/2$};
  \node[font=\scriptsize] at (8.00,3.73) {$\widetilde\theta$};

  \draw[seammap]
    (13.15,3.15)--(13.54,3.15)--(13.54,.72)--(7.61,.72)--(7.61,1.80)--(8.00,1.80);
  \node[fill=white,inner sep=3pt,font=\scriptsize,text=figSeam]
    at (10.57,.72) {rotated seam};
  \draw[seamline,fill=white] (8.00,1.80) circle[radius=2.0pt];
  \draw[pathline,arrow] (8.00,3.15)--(13.15,3.15);
  \fill[figRotation] (8.00,3.15) circle[radius=1.7pt];
  \fill[figRotation] (13.15,3.15) circle[radius=1.7pt];
  \node[text=figRotation] at (10.57,2.72) {$\widetilde\theta=0$};
  \node[font=\scriptsize,text=black!65] at (10.57,2.27)
    {$\widetilde\theta=\theta+\rho\tau$};

  % The torsion coordinate and endpoint identifications.
  \foreach \a/\b/\c in {.90/6.05/3.47,8.00/13.15/10.57} {
    \node[font=\scriptsize] at (\a,1.19) {$0$};
    \node[font=\scriptsize] at (\b,1.19) {$2\pi$};
    \draw[arrow,figTorsion] ({\c-.42},1.19)--({\c+.32},1.19);
    \node[anchor=west,inner sep=2pt,text=figTorsion,font=\scriptsize]
      at ({\c+.32},1.19) {$\tau$};
  }
  \node[font=\scriptsize] at (3.47,.19)
    {$(\theta,2\pi)\sim(\theta,0)$};
  \node[font=\scriptsize] at (10.57,.19)
    {$(\widetilde\theta,2\pi)\sim(\widetilde\theta-2\pi\rho,0)$};
\end{tikzpicture}
\caption{One torsional loop in two frames for the same illustrative
two-rotor model. Solid traces show zero-angular-momentum transport,
and the dials show the same initial and final laboratory orientations.
The periodic frame carries the rotation in its connection.
The parallel frame carries it in the endpoint identification.
In both descriptions the accumulated rotation is $-2\pi\rho=-\pi/2$.}
\label{fig:connection-and-seam}
\end{figure}
